\chapter{Kajian Pustaka}

\section{Teori Dasar}
\subsection{Process Mining pada Domain Pemeliharaan}
\textit{Process mining} merupakan pendekatan analitik yang mengekstraksi pengetahuan dari \textit{event logs} untuk memahami hubungan antar-proses dalam suatu sistem~\cite{horn2019application}. Dalam domain pemeliharaan, pendekatan ini memungkinkan analisis tidak hanya pada tingkat peralatan secara individu, tetapi juga mencakup keterkaitan antar-unit secara lebih menyeluruh. Dengan merekam \textit{traces} aktivitas yang terjadi secara berurutan, \textit{process mining} dapat mengungkap pola perilaku mesin, ketergantungan antar-proses, serta hambatan dalam alur produksi yang sebelumnya sulit dideteksi~\cite{horn2019application}. Hasil analisis ini memberikan dasar bagi pengambilan keputusan yang lebih tepat dalam upaya optimalisasi peralatan dan perpanjangan masa pakai mesin~\cite{horn2019application}.

\subsection{Process Mining dan Process Discovery}
\textit{Process mining} merupakan disiplin ilmu yang menjembatani \textit{data science} dengan manajemen proses bisnis (\textit{Business Process Management}) untuk mengekstraksi wawasan dari \textit{event logs} \cite{vanderAalst2016_process_mining}. Salah satu kapabilitas utamanya adalah \textit{process discovery}, yaitu teknik untuk memodelkan proses bisnis secara otomatis tanpa informasi apriori mengenai bagaimana proses seharusnya berjalan \cite{augusto2019automated}.

Meskipun algoritma \textit{process discovery} konvensional dapat menghasilkan model dari data terstruktur, penerapannya pada lingkungan kompleks seperti MRO menghadapi tantangan terkait tingginya \textit{noise} dan variabilitas perilaku. Hal ini menyebabkan model yang dihasilkan menjadi rumit dan sulit dipahami (\textit{spaghetti models}) \cite{augusto2019automated}. Selain itu, algoritma tradisional kerap mengabaikan pengetahuan domain yang tersirat dalam dokumen teks, sehingga model yang dihasilkan mungkin akurat secara statistik tetapi tidak sepenuhnya patuh pada logika bisnis atau regulasi keselamatan \cite{norouzifar2024}.

\subsection{Generative Artificial Intelligence dan Large Language Models}
Generative AI dan Large Language Model (LLM) merupakan dua konsep yang saling terkait erat dan bersama-sama mendorong revolusi dalam pemrosesan bahasa alami serta berbagai aplikasi kecerdasan buatan modern. Generative AI adalah payung besar untuk teknologi yang mampu menciptakan konten baru---mulai dari teks, gambar, hingga audio---dengan mempelajari pola dari data yang sangat besar. Salah satu pencapaian paling menonjol dari generative AI adalah kemunculan LLM, seperti GPT, BERT, dan PaLM2, yang secara khusus dirancang untuk memahami, memproses, dan menghasilkan bahasa alami dengan tingkat kemiripan yang sangat tinggi dengan manusia \cite{hagos2024recentLLM}.

LLM dibangun di atas arsitektur transformer dan dilatih menggunakan korpus teks dalam skala masif, sehingga mampu mengenali pola bahasa yang kompleks dan menghasilkan respons yang relevan untuk berbagai tugas, seperti tanya jawab, rangkuman, terjemahan, hingga pembuatan kode \cite{yan2024generativeLLM}.

\subsection{PM4Py (Process Mining for Python)}
PM4Py adalah pustaka perangkat lunak \textit{open-source} yang menyediakan berbagai algoritma \textit{process mining} modern dalam ekosistem Python \cite{berti2019pm4py}. Pustaka ini mendukung algoritma \textit{Inductive Miner}, \textit{Heuristic Miner}, \textit{conformance checking}, dan analisis kinerja. Fleksibilitasnya memungkinkan integrasi keluaran anotasi dari LLM, seperti label \texttt{action\_type}, ke dalam alur kerja \textit{process discovery} sehingga ideal untuk pendekatan hibrida \textit{AI-assisted process mining}.

\subsection{Petri Net}
Sebuah \textit{Petri net} dapat dinyatakan sebagai tripel $N = (P, T, F)$ dengan $P$ himpunan \textit{place}, $T$ himpunan \textit{transition}, dan $F \subseteq (P \times T) \cup (T \times P)$ relasi alir yang menghubungkan keduanya. Kondisi proses direpresentasikan oleh distribusi \textit{token} pada \textit{place}, yang disebut \textit{marking}. Sebuah \textit{transition} dapat menyala (\textit{fire}) bila seluruh \textit{place} masukannya memiliki \textit{token}, dan penyalaan itu memindahkan \textit{token} sehingga \textit{marking} berubah. Notasi ini menjadi keluaran utama tahap \textit{process discovery} pada Tugas Akhir, sebab algoritma \textit{Inductive Miner} menghasilkan model dalam bentuk \textit{Petri net} yang \textit{sound} dan \textit{block-structured} \cite{vanderAalst2016_process_mining}.

\subsection{Directly-Follows Graph (DFG)}
\textit{Directly-Follows Graph} (DFG) adalah graf berarah yang merangkum relasi urutan langsung antaraktivitas pada \textit{event log}. Setiap simpul mewakili satu aktivitas, sedangkan sebuah busur $a \rightarrow b$ dibentuk bila aktivitas $a$ langsung diikuti aktivitas $b$ dalam suatu \textit{trace}, umumnya disertai bobot frekuensi kemunculan. DFG berperan sebagai struktur antara yang dibaca algoritma \textit{process discovery} untuk menyimpulkan pola proses, namun representasi ini rawan terhadap \textit{noise} ketika log memuat banyak variasi label \cite{augusto2019automated}.

\subsection{Inductive Miner dan Inductive Miner infrequent (IMf)}
\textit{Inductive Miner} adalah keluarga algoritma \textit{process discovery} yang membangun model dengan membelah DFG secara rekursif menjadi \textit{process tree}. Pada tiap langkah, algoritma mendeteksi pola pembelahan (\textit{cut}) berupa pilihan eksklusif, urutan, paralel, atau pengulangan, lalu memetakan hasilnya menjadi \textit{Petri net} yang dijamin \textit{sound}. Varian \textit{Inductive Miner infrequent} (IMf) menambahkan \textit{noise filtering} dengan membuang busur DFG yang frekuensinya berada di bawah \textit{noise threshold} tertentu. Nilai \textit{noise threshold} 0 mempertahankan seluruh perilaku, sedangkan nilai yang lebih tinggi menghasilkan model lebih ringkas dengan mengabaikan jalur jarang. Penelitian ini memakai IMf dengan \textit{noise threshold} 0,2 \cite{leemans2013discovering, augusto2019automated}.

\subsection{Event Abstraction}
\textit{Event abstraction} adalah proses menyatukan \textit{events} bertingkat rinci menjadi aktivitas bertingkat lebih tinggi yang lebih bermakna, sehingga model proses yang ditemukan menjadi lebih ringkas dan mudah ditafsirkan \cite{tax2016event}. Pada penelitian ini, abstraksi semantik dimaksudkan sebagai bentuk \textit{event abstraction} yang memetakan deskripsi aktivitas mentah ke kategori taksonomi melalui LLM, sebelum model proses dibangun.

\section{Conformance Checking}
\textit{Conformance checking} merupakan tahap evaluasi yang mengukur kesesuaian antara model proses yang ditemukan dengan perilaku yang terekam dalam \textit{event log}~\cite{carmona2018conformance}. Empat metrik utama digunakan dalam evaluasi ini, yaitu \textit{fitness}, \textit{precision}, \textit{simplicity}, dan \textit{generalization}, yang dirumuskan sebagai kerangka kualitas model proses oleh Buijs dkk.~\cite{buijs2014quality}.

\subsection{Fitness}
\textit{Fitness} merupakan dimensi paling fundamental yang mengukur sejauh mana model proses yang ditemukan mampu mereproduksi \textit{traces} aktivitas yang terekam dalam \textit{event log}. Sebuah model dikatakan memiliki \textit{fitness} sempurna (nilai $1.0$) apabila seluruh \textit{traces} dalam log dapat di-\textit{replay} tanpa penyimpangan~\cite{vanderAalst2016_process_mining}. Secara kuantitatif, nilai \textit{fitness} dihitung menggunakan pendekatan \textit{alignment-based fitness} yang berakar pada analisis kesesuaian berbasis biaya oleh Adriansyah dkk.~\cite{adriansyah2011conformance, adriansyah2012replaying} sebagai berikut:

\begin{equation}
    Fitness(L, M) = 1 - \frac{\sum_{\sigma \in L} n(\sigma) \times \delta_{opt}(\sigma, M)}{\sum_{\sigma \in L} n(\sigma) \times c_{wc}(\sigma)}
    \label{eq:fitness_formula}
\end{equation}

Di mana $n(\sigma)$ adalah frekuensi kemunculan suatu \textit{trace} dalam log, $\delta_{opt}$ merupakan biaya penyelarasan optimal (\textit{alignment cost}), dan $c_{wc}$ adalah biaya terburuk (\textit{worst-case cost})~\cite{adriansyah2012replaying}.

\subsection{Precision}
Berbeda dengan \textit{fitness}, \textit{precision} membalik sudut pandang evaluasi dengan menanyakan sejauh mana perilaku yang dimungkinkan oleh model benar-benar tercatat dalam \textit{event log}~\cite{carmona2018conformance, carmona2022conformance}. Sebuah model dikatakan memiliki \textit{precision} rendah apabila terlalu permisif, yaitu mengizinkan banyak perilaku yang tidak pernah terjadi dalam data aktual~\cite{carmona2018conformance}. Secara kuantitatif, \textit{precision} dihitung menggunakan pendekatan \textit{escaping edges} yang pertama kali dirumuskan Muñoz-Gama dan Carmona~\cite{munozgama2010precision, vanderAalst2016_process_mining} sebagai berikut:

\begin{equation}
    Precision(L, M) = \frac{1}{|E|} \sum_{e \in E} \frac{|a_L(e)|}{|a_M(e)|}
    \label{eq:precision_formula}
\end{equation}

Di mana $E$ adalah himpunan \textit{state} yang dikunjungi, $a_{L}(e)$ merupakan jumlah aktivitas yang terekam dalam log pada \textit{state} tersebut, dan $a_{M}(e)$ adalah jumlah aktivitas yang diizinkan oleh model pada \textit{state} yang sama.

\subsection{Simplicity}
\textit{Simplicity} mengukur sejauh mana struktur model proses tetap sederhana dan tidak berlebihan dalam jumlah elemen. Salah satu rumusan yang lazim dipakai dan diterapkan pada PM4Py adalah pendekatan \textit{inverse arc degree}, yang menilai kesederhanaan dari rata-rata derajat keterhubungan simpul \textit{Petri net}~\cite{buijs2014quality, berti2019pm4py}. Bila $\bar{d}$ menyatakan rata-rata derajat busur, yaitu jumlah busur masuk dan keluar seluruh \textit{place} pada model, dan $k$ adalah derajat acuan yang lazimnya bernilai 2, maka \textit{simplicity} dirumuskan sebagai berikut.

\begin{equation}
    Simplicity(M) = \frac{1}{1 + \log(\bar{d}/k)}
    \label{eq:simplicity_formula}
\end{equation}

Nilai mendekati 1 menandakan model sederhana dengan keterhubungan wajar, sedangkan nilai yang menurun menandakan struktur yang makin padat. Model dengan \textit{simplicity} rendah sering disebut sebagai \textit{spaghetti model}, yaitu model yang memiliki struktur sangat ruwet dan sulit dipahami oleh manusia~\cite{vanderAalst2016_process_mining, carmona2018conformance}.

\subsection{Generalization}
\textit{Generalization} menilai sejauh mana model mampu menangkap perilaku wajar yang mungkin belum terekam pada log, sehingga model tidak \textit{overfitting} terhadap data pelatihan. Buijs dkk.~\cite{buijs2014quality} merumuskan \textit{generalization} dengan memanfaatkan frekuensi kunjungan tiap \textit{transition}. Bila $T_M$ himpunan \textit{transition} pada model dan $n_t$ banyaknya eksekusi \textit{transition} $t$ selama \textit{replay} pada log $L$, maka \textit{generalization} dirumuskan sebagai berikut.

\begin{equation}
    Generalization(L, M) = \prod_{t \in T_M} \left(1 - \frac{1}{n_t}\right)
    \label{eq:generalization_formula}
\end{equation}

\textit{Transition} yang sering dieksekusi memberi kontribusi penalti kecil sehingga model dinilai lebih tergeneralisasi, sedangkan \textit{transition} yang jarang dikunjungi menurunkan nilai \textit{generalization}. Pendekatan alternatif berbasis \textit{anti-alignment} yang menyatukan pengukuran \textit{precision} dan \textit{generalization} dikemukakan van Dongen dkk.~\cite{vandongen2016unified}, dan dapat dirujuk sebagai pembanding.

\section{Penelitian Terkait}
Bagian ini mengulas penelitian terdahulu yang menjadi pijakan metode Tugas Akhir, khususnya mengenai cara kecerdasan buatan bekerja sama dengan \textit{process mining} untuk memperjelas struktur proses. Penekanannya bukan pada penyaringan data, melainkan pada bagaimana makna aktivitas diolah agar model proses yang dihasilkan lebih ringkas dan mudah ditafsirkan.

\subsection{LLM untuk Menjembatani Pengetahuan Domain dan Process Discovery}
Norouzifar dkk. menyoroti satu kendala klasik dalam \textit{process discovery} otomatis, yaitu algoritma sering kesulitan memisahkan variasi kerja yang wajar dari catatan yang keliru. Mereka menempatkan LLM sebagai penerjemah yang membaca deskripsi tekstual aturan, lalu mengubahnya menjadi batasan deklaratif sederhana, misalnya keharusan satu aktivitas mendahului aktivitas lain. Batasan itu kemudian memandu algoritma \textit{Inductive Miner} agar membuang jalur proses yang tidak koheren, sehingga model akhir lebih patuh pada logika domain dibandingkan jika hanya bersandar pada \textit{raw data} \cite{norouzifar2024}.

Pendekatan ini menjadi rujukan langsung karena memakai algoritma penemuan model yang sama dengan Tugas Akhir. Perbedaannya terletak pada peran LLM. Norouzifar dkk. memanfaatkan LLM untuk mengekstrak \textit{constraint} proses, sementara penelitian ini memakai LLM untuk memberi label kategori tindakan pada tiap \textit{event}. Celah yang tersisa jelas, yaitu penyaringan berbasis aturan belum tentu menjawab persoalan utama pada log pemeliharaan, yang justru bersumber dari keberagaman penulisan label aktivitas.

\subsection{Event Abstraction Berbasis Supervised Learning}
Tax dkk. menelaah persoalan ketika \textit{events} pada log terlalu rinci sehingga algoritma penemuan menghasilkan model yang kusut dan sukar dibaca. Mereka menunjukkan bahwa struktur proses masih dapat ditemukan apabila log lebih dahulu diabstraksikan ke tingkat granularitas yang lebih tinggi. Untuk itu, mereka memakai teknik \textit{supervised learning}, khususnya \textit{Linear-chain Conditional Random Fields}, dengan memanfaatkan anotasi bermakna pada sebagian \textit{trace} sebagai data latih, lalu merepresentasikan tiap \textit{event} sebagai vektor fitur. Pengujian pada log nyata sebuah lingkungan \textit{smart home} memperlihatkan model proses yang jauh lebih terstruktur setelah abstraksi \cite{tax2016event}.

Gagasan inilah yang menjadi akar konseptual penelitian ini, sebab abstraksi semantik pada dasarnya merupakan bentuk \textit{event abstraction}. Letak pembedanya ada pada cara abstraksi dilakukan. Tax dkk. menuntut data anotasi terlatih dan model probabilistik klasik, sedangkan Tugas Akhir ini menyandarkan abstraksi semantik pada LLM dengan dua belas kategori taksonomi sebagai batas, tanpa pelatihan ulang. Hal yang belum dijawab oleh Tax dkk. adalah bagaimana abstraksi dapat berjalan ketika label latih tidak tersedia, kondisi yang lazim ditemui pada catatan pemeliharaan lapangan.

\subsection{LLM untuk Tugas Process Mining yang Peka Makna}
Rebmann dkk. menguji secara sistematis sejauh mana model bahasa mampu menyelesaikan tugas \textit{process mining} yang menuntut pemahaman makna aktivitas. Mereka mendefinisikan lima tugas, mencakup dua bentuk deteksi anomali semantik, prediksi aktivitas berikutnya secara semantik, dan dua varian \textit{process discovery} semantik, lalu membandingkan kinerja LLM pada skenario \textit{in-context learning} dengan skenario \textit{fine-tuning}. Temuan utamanya, model bahasa cenderung lemah ketika dipakai apa adanya atau dengan contoh minimal, namun kinerjanya meningkat tajam saat disesuaikan melalui \textit{fine-tuning} lintas beragam jenis proses dan industri \cite{rebmann2025semantics}.

Studi ini menegaskan premis dasar Tugas Akhir, yaitu makna aktivitas memang dapat dieksploitasi LLM untuk kepentingan penemuan proses. Perbedaannya bersifat penerapan. Rebmann dkk. menilai kapabilitas umum LLM lintas \textit{benchmark} dan banyak bersandar pada \textit{fine-tuning}, sementara penelitian ini menerapkan LLM secara langsung untuk abstraksi semantik label aktivitas pada satu domain pemeliharaan pembangkit listrik, lalu mengukur dampaknya terhadap kualitas model melalui perbandingan \textit{as-is} dan \textit{enriched}. Ruang yang belum tergarap adalah penerapan terfokus pada data operasional nyata dengan evaluasi \textit{conformance} yang utuh.

\section{Perbandingan Penelitian}
Berdasarkan tinjauan tersebut, Tabel~\ref{tab:perbandingan-terkait} merangkum posisi Tugas Akhir ini dibandingkan ketiga referensi utama.

\begin{table}[htbp]
    \centering
    \small
    \renewcommand{\arraystretch}{1.4}
    \begin{tabularx}{\textwidth}{|p{2.4cm}|X|X|X|X|}
        \hline
        \textbf{Aspek} &
        \textbf{Norouzifar dkk.~\cite{norouzifar2024}} &
        \textbf{Tax dkk.~\cite{tax2016event}} &
        \textbf{Rebmann dkk.~\cite{rebmann2025semantics}} &
        \textbf{Penelitian Ini (TA)} \\
        \hline
        \textbf{Fokus Utama} &
        Penyaringan jalur proses lewat batasan deklaratif hasil ekstraksi LLM &
        \textit{Event abstraction} dari tingkat rinci ke tingkat bermakna agar model lebih terstruktur &
        Evaluasi kapabilitas LLM pada lima tugas \textit{process mining} yang bergantung pada makna &
        Abstraksi semantik berbasis taksonomi via LLM untuk \textit{process discovery} \\
        \hline
        \textbf{Metode AI} &
        LLM mengekstrak \textit{constraint} dari teks secara \textit{zero-shot} &
        \textit{Supervised learning}, \textit{Linear-chain Conditional Random Fields}, bukan LLM &
        LLM melalui \textit{in-context learning} dan \textit{fine-tuning} &
        LLM (\texttt{gemini-3.1-flash-lite}) untuk pelabelan kategori dengan batas taksonomi, tanpa \textit{fine-tuning} \\
        \hline
        \textbf{Sumber Pengetahuan} &
        Deskripsi aturan tekstual dan pengetahuan internal LLM &
        Anotasi bermakna pada sebagian \textit{trace} sebagai data latih &
        Pengetahuan terenkode LLM dan \textit{dataset benchmark} &
        Konteks \textit{event} dari WBS dan \textit{job card}, taksonomi dua belas kategori, dan pengetahuan LLM \\
        \hline
        \textbf{Domain dan Data} &
        Proses bisnis umum &
        Lingkungan \textit{smart home} berbasis sensor &
        Lintas industri melalui \textit{dataset benchmark} &
        Pemeliharaan pembangkit listrik (PLTD), 256 baris pekerjaan \\
        \hline
        \textbf{Kontribusi dan Perbedaan} &
        Model lebih patuh logika domain, namun bertumpu pada ekstraksi aturan &
        Model lebih terstruktur dari log granular, namun membutuhkan label latih &
        Membuktikan LLM unggul saat \textit{fine-tuned}, namun berhenti pada evaluasi kapabilitas &
        Reduksi aktivitas unik dan peningkatan \textit{generalization} melalui abstraksi semantik, disertai komparasi \textit{as-is} dan \textit{enriched} dengan \textit{conformance} \\
        \hline
    \end{tabularx}
    \caption{Perbandingan Penelitian Terkait}
    \label{tab:perbandingan-terkait}
\end{table}

Benang merah dari ketiga studi memperlihatkan satu celah yang konsisten. Norouzifar dkk. memakai LLM untuk menyusun \textit{constraint}, Tax dkk. mengandalkan \textit{supervised learning} klasik yang menuntut data latih, dan Rebmann dkk. berhenti pada pengukuran kapabilitas lintas \textit{benchmark}. Belum ada yang memadukan abstraksi semantik berbasis taksonomi melalui LLM dengan komparasi model \textit{as-is} dan \textit{enriched} serta pemeriksaan \textit{conformance} pada log pemeliharaan pembangkit listrik. Ruang itulah yang ditempati Tugas Akhir ini.

\newpage
